% ============================================================
% The Collision Periodic Table
% ============================================================

\documentclass[12pt]{amsart}

\usepackage{amsmath,amssymb,amsthm}
\usepackage{booktabs}
\usepackage{microtype}
\usepackage{url}

\newtheorem{theorem}{Theorem}
\newtheorem{lemma}[theorem]{Lemma}
\newtheorem{proposition}[theorem]{Proposition}
\newtheorem{corollary}[theorem]{Corollary}
\theoremstyle{definition}
\newtheorem{definition}[theorem]{Definition}
\theoremstyle{remark}
\newtheorem{remark}[theorem]{Remark}
\newtheorem{observation}[theorem]{Observation}

\title{The Collision Periodic Table}

\author{Alexander S.\ Petty}

\email{alexander.petty@gmail.com}
\date{December 2023; updated April 2026}
\subjclass[2020]{11A63, 11N13, 11A07}

\begin{document}
\begin{abstract}
For a prime $p > b^2$ not dividing the base~$b$, the
lag-$1$ signed collision count
\[
S(p) = C(b \bmod p) - \left\lfloor \frac{p-1}{b}\right\rfloor
\]
is an integer determined by $p \bmod b^2$. The
$\varphi(b^2)$ values form the \emph{collision periodic
table}, a fixed invariant of the base.

The table satisfies the antisymmetry law
$S(a) + S(b^2 - a) = -1$ for every coprime class~$a$,
the range is $[-(b{-}1),\; b{-}2]$, the mean is exactly
$-1/2$, and exactly half the classes are negative. These
properties are proved by elementary complement counting
and floor arithmetic. A brief final note records the later
analytic interpretation developed in 2026.
\end{abstract}
\maketitle

% ============================================================
\section{Introduction}

This paper exhibits the lag-$1$ collision periodic table:
a finite, computable table that assigns an integer to each
coprime residue class modulo~$b^2$ and determines the
collision behavior of every sufficiently large prime.

The point is elementary. At lag~$1$, the specific
collision count at multiplier~$b$ can be written as a
finite slice count depending only on the residue class of
$p$ modulo~$b^2$. This produces a fixed table with
$\varphi(b^2)$ cells, one integer per coprime class.

The resulting table is not approximate. It is exact: an
integer per class, constant for all primes in that class,
with proved structural properties. In base~$10$, the table
has $40$ entries indexed by the coprime classes
modulo~$100$. The last two decimal digits of a prime
determine its collision fingerprint.

% ============================================================
\section{The Fingerprint Theorem}

\begin{theorem}[Finite determination at lag~$1$]
\label{thm:fingerprint}
For any base $b \ge 2$ and prime $p > b^2$ with
$p \nmid b$, the signed collision count
\[
S(p) = C(b \bmod p) -
\left\lfloor \frac{p-1}{b} \right\rfloor
\]
is an integer that depends only on $p \bmod b^2$.
\end{theorem}

\begin{proof}
Set $m=b^2$. For each $r \in \{1,\dots,p-1\}$, define the
slice index
\[
n(r)=\left\lfloor \frac{mr}{p}\right\rfloor
\in \{0,\dots,m-1\}.
\]
Since the slice width $p/m>1$, each slice contains
integers. On the slice with index $n$, both the digit
function and the digit of the image under multiplication
by~$b$ are constant:
\[
\left\lfloor \frac{br}{p}\right\rfloor
= \left\lfloor \frac{n}{b}\right\rfloor,
\qquad
\delta(br \bmod p)= n \bmod b.
\]
Thus a slice contributes to the collision count exactly
when
\[
\left\lfloor \frac{n}{b}\right\rfloor = n \bmod b.
\]
The diagonal slices are therefore
\begin{align*}
G &= \bigl\{n \in \{0,\dots,b^2{-}1\} :
\lfloor n/b\rfloor = n \bmod b\bigr\} \\
&= \{r(b{+}1) : 0 \le r \le b{-}1\}.
\end{align*}
Hence
\[
C(b \bmod p)
= -1 + \sum_{n \in G}
\left(
\left\lfloor \frac{(n+1)p}{b^2}\right\rfloor
- \left\lfloor \frac{np}{b^2}\right\rfloor
\right),
\]
where the $-1$ removes $r=0$ from the $n=0$ slice.

Write $p=b^2Q_2+a$ with $1\le a<b^2$ and $\gcd(a,b)=1$.
Then
\[
\left\lfloor \frac{(n+1)p}{b^2}\right\rfloor
- \left\lfloor \frac{np}{b^2}\right\rfloor
= Q_2 + \left\lfloor \frac{(n+1)a}{b^2}\right\rfloor
- \left\lfloor \frac{na}{b^2}\right\rfloor.
\]
Since
\[
\left\lfloor \frac{p-1}{b}\right\rfloor
= bQ_2 + \left\lfloor \frac{a-1}{b}\right\rfloor,
\]
it follows that
\[
S(p)= -1 - \left\lfloor \frac{a-1}{b}\right\rfloor
+ \sum_{n\in G}
\left(
\left\lfloor \frac{(n+1)a}{b^2}\right\rfloor
- \left\lfloor \frac{na}{b^2}\right\rfloor
\right),
\]
which depends only on $a=p \bmod b^2$.
\end{proof}

\begin{definition}
The \emph{collision periodic table} in base $b$ is the function
$\mathcal{T} \colon (\mathbb{Z}/b^2\mathbb{Z})^{*}
\to \mathbb{Z}$ defined by $\mathcal{T}(a) = S(p)$ for
any prime $p \equiv a \pmod{b^2}$ with $p > b^2$.
\end{definition}

% ============================================================
\section{The Antisymmetry Law}

\begin{lemma}[Lag-$1$ closed form]\label{lem:closed}
Let $p>b^2$ be prime with $p\nmid b$, and write
$p=b^2Q_2+a$ with $1\le a<b^2$ and $\gcd(a,b)=1$.
Then
\[
C(b,p)=bQ_2+f(a)-1,
\]
where
\[
f(a)=\#\{d\in\{0,\dots,b-1\} :
(b+1)da \bmod b^2 \ge b^2-a\}.
\]
\end{lemma}

\begin{proof}
From Theorem~\ref{thm:fingerprint},
\[
C(b \bmod p)
= -1 + \sum_{n \in G}
\left(
\left\lfloor \frac{(n+1)p}{b^2}\right\rfloor
- \left\lfloor \frac{np}{b^2}\right\rfloor
\right),
\]
where $G=\{d(b+1):0\le d\le b-1\}$. Writing
$n=d(b+1)$ and $p=b^2Q_2+a$ gives
\[
\left\lfloor \frac{(n+1)p}{b^2}\right\rfloor
- \left\lfloor \frac{np}{b^2}\right\rfloor
= Q_2 +
\left\lfloor \frac{(d(b+1)+1)a}{b^2}\right\rfloor
- \left\lfloor \frac{d(b+1)a}{b^2}\right\rfloor.
\]
Since $0<a<b^2$, the final difference is either $0$ or $1$.
It equals $1$ exactly when the fractional part of
$d(b+1)a/b^2$ is at least $1-a/b^2$, equivalently when
\[
(b+1)da \bmod b^2 \ge b^2-a.
\]
Thus each diagonal slice contributes $Q_2$ plus the
indicator of that threshold event. Summing over
$d=0,\dots,b-1$ gives the formula.
\end{proof}

\begin{theorem}[Antisymmetry]\label{thm:antisym}
For every base $b \ge 3$ and every coprime class
$a \bmod b^2$:
\[
\mathcal{T}(a) + \mathcal{T}(b^2 - a) = -1.
\]
\end{theorem}

\begin{proof}
By Lemma~\ref{lem:closed},
the collision count at multiplier $g=b$ satisfies
$Q = \lfloor (p{-}1)/b \rfloor = b Q_2 +
\lfloor (a{-}1)/b \rfloor$:
\[
\mathcal{T}(a) = f(a) - 1 - \left\lfloor
\frac{a-1}{b} \right\rfloor.
\]

Three identities close the proof.

\medskip\noindent
\emph{(i) Complementary counting.}
Write $r_d = (b{+}1) d \cdot a \bmod b^2$. For $d \ge 1$,
$r_d \ne 0$ (since $\gcd(a, b^2) = 1$ and
$\gcd(b{+}1, b^2) = 1$). The complement of $r_d$ is
$b^2 - r_d = (b{+}1) d \cdot (b^2{-}a) \bmod b^2$.
The condition $r_d \ge b^2 - a$ for $f(a)$ becomes
$b^2 - r_d \le a$, which is the complement of the
condition $b^2 - r_d \ge a$ for $f(b^2{-}a)$.
For $d \ge 1$: every $d$ is counted by exactly one of
$f(a)$ or $f(b^2{-}a)$, except when $r_d = b^2 - a$
(counted by both). Including $d = 0$ (where
$r_0 = 0$, counted by neither):
$f(a) + f(b^2{-}a) = b$.

\medskip\noindent
\emph{(ii) Floor sum.}
$(a{-}1) + (b^2{-}a{-}1) = b^2 - 2$. The remainders
modulo~$b$ are $R = (a{-}1) \bmod b$ and
$R' = (b^2{-}a{-}1) \bmod b = b{-}2{-}R$. Since
$R + R' = b - 2 < b$:
$\lfloor (a{-}1)/b \rfloor +
\lfloor (b^2{-}a{-}1)/b \rfloor = (b^2 - 2 - (b-2))/b
= b - 1$.

\medskip\noindent
\emph{(iii) Assembly.}
$\mathcal{T}(a) + \mathcal{T}(b^2{-}a) =
[f(a) + f(b^2{-}a)] - 2 -
[\lfloor(a{-}1)/b\rfloor + \lfloor(b^2{-}a{-}1)/b\rfloor]
= b - 2 - (b - 1) = -1$.
\end{proof}

\begin{corollary}\label{cor:range}
For every base $b \ge 3$:
\[
\mathcal{T}(b{-}1) = b - 2, \qquad
\mathcal{T}(b^2 - b + 1) = -(b-1).
\]
The range of $\mathcal{T}$ is contained in
$[-(b{-}1),\; b{-}2]$.
\end{corollary}

\begin{proof}
For $a=b-1$, the formula from the antisymmetry proof gives
\[
\mathcal T(a)=f(a)-1-\left\lfloor \frac{a-1}{b}\right\rfloor
=f(b-1)-1.
\]
Now
\[
r_d=(b+1)d(b-1)\equiv -d \pmod{b^2},
\]
so for each $d=1,\dots,b-1$ one has
$r_d=b^2-d \ge b^2-b+1 = b^2-a$, while $d=0$ is not
counted. Thus $f(b-1)=b-1$, hence
$\mathcal T(b-1)=b-2$.

The complementary class $b^2-b+1$ then satisfies
\[
\mathcal T(b^2-b+1)=-1-\mathcal T(b-1)=-(b-1)
\]
by antisymmetry.

Finally, $d=0$ is never counted in $f(a)$, so
$0\le f(a)\le b-1$ for every coprime class~$a$. Hence
\[
\mathcal T(a)=f(a)-1-\left\lfloor \frac{a-1}{b}\right\rfloor
\le (b-1)-1=b-2.
\]
Applying antisymmetry gives
$\mathcal T(a)\ge -(b-1)$.
\end{proof}

% ============================================================
\section{The Table}

\subsection*{Base 10}

The $40$ coprime classes modulo~$100$ and their collision
fingerprints at lag~$1$:

\begin{table}[h]
\centering
\begin{tabular}{r*{10}{r}}
\toprule
& 0 & 1 & 2 & 3 & 4 & 5 & 6 & 7 & 8 & 9 \\
\midrule
0 & . & $0$ & . & $+2$ & . & . & . & $0$ & . & $+8$ \\
1 & . & $-1$ & . & $-1$ & . & . & . & $+1$ & . & $-1$ \\
2 & . & $0$ & . & $-2$ & . & . & . & $+6$ & . & $0$ \\
3 & . & $-1$ & . & $-1$ & . & . & . & $-3$ & . & $-1$ \\
4 & . & $-4$ & . & $0$ & . & . & . & $-2$ & . & $0$ \\
5 & . & $-1$ & . & $+1$ & . & . & . & $-1$ & . & $+3$ \\
6 & . & $0$ & . & $+2$ & . & . & . & $0$ & . & $0$ \\
7 & . & $-1$ & . & $-7$ & . & . & . & $+1$ & . & $-1$ \\
8 & . & $0$ & . & $-2$ & . & . & . & $0$ & . & $0$ \\
9 & . & $-9$ & . & $-1$ & . & . & . & $-3$ & . & $-1$ \\
\bottomrule
\end{tabular}
\caption{Collision periodic table in base~$10$. Row is the
tens digit, column is the units digit. Dots indicate classes
not coprime to~$100$. The entry is
$S = C(10 \bmod p) - \lfloor(p{-}1)/10\rfloor$, constant
for all primes $p > 100$ in that class.}
\end{table}

The table has $8$ positive entries, $12$ zeros, and $20$
negative entries. The extremes are $\mathcal{T}(9) = +8$ and
$\mathcal{T}(91) = -9$, confirming Corollary~\ref{cor:range}.

\subsection*{Universality across bases}

The structural properties hold for every base tested:

\begin{table}[h]
\centering
\begin{tabular}{rrrrrrrl}
\toprule
$b$ & $\varphi(b^2)$ & $S_{\min}$ & $S_{\max}$ &
$+$ & $0$ & $-$ & Antisymmetry \\
\midrule
$3$  &  $6$  & $-2$ & $+1$ &  1 &  2 &  3 & $S(a)+S({-}a)=-1$ \\
$5$  & $20$  & $-4$ & $+3$ &  4 &  6 & 10 & $S(a)+S({-}a)=-1$ \\
$6$  & $12$  & $-5$ & $+4$ &  2 &  4 &  6 & $S(a)+S({-}a)=-1$ \\
$7$  & $42$  & $-6$ & $+5$ & 10 & 11 & 21 & $S(a)+S({-}a)=-1$ \\
$8$  & $32$  & $-7$ & $+6$ &  7 &  9 & 16 & $S(a)+S({-}a)=-1$ \\
$10$ & $40$  & $-9$ & $+8$ &  8 & 12 & 20 & $S(a)+S({-}a)=-1$ \\
$12$ & $48$  &$-11$ &$+10$ & 11 & 13 & 24 & $S(a)+S({-}a)=-1$ \\
\bottomrule
\end{tabular}
\caption{Summary across bases. In every case:
$S_{\min} = -(b{-}1)$, $S_{\max} = b - 2$,
antisymmetry constant $= -1$, and negative classes
outnumber positive.}
\end{table}

% ============================================================
\section{The Sign Balance}

\begin{corollary}[Negative majority]\label{cor:negative}
Exactly $\varphi(b^2)/2$ classes have $\mathcal{T}(a) < 0$.
\end{corollary}

\begin{proof}
The antisymmetry pairs each class $a$ with $-a$, giving
$\varphi(b^2)/2$ pairs. In each pair,
$\mathcal{T}(a) + \mathcal{T}(-a) = -1 < 0$, so at least
one member is negative. If $\mathcal{T}(a) = 0$ then
$\mathcal{T}(-a) = -1 < 0$. If $\mathcal{T}(a) > 0$ then
$\mathcal{T}(-a) = -1 - \mathcal{T}(a) < -1 < 0$. In every
case, exactly one member per pair is negative.
\end{proof}

\begin{corollary}[Mean]\label{cor:mean}
\[
\frac{1}{\varphi(b^2)} \sum_{\substack{a \bmod b^2 \\
\gcd(a,b)=1}} \mathcal{T}(a) = -\frac{1}{2}.
\]
\end{corollary}

\begin{proof}
Each pair $(a, -a)$ contributes $\mathcal{T}(a) +
\mathcal{T}(-a) = -1$. There are $\varphi(b^2)/2$ pairs.
The total is $-\varphi(b^2)/2$.
\end{proof}

% ============================================================
\section*{Later Addition (2026)}

This note was written to record the elementary lag-$1$
table and its antisymmetry. Later work reinterpreted the
same table through Fourier analysis on
$(\mathbb Z/b^2\mathbb Z)^{\times}$. In that later
framework, the table coefficients are supported on the
primitive odd sector and factor through Bernoulli numbers
and diagonal character sums. That analytic story belongs
to the later transform and spectrum papers, not to the
2023 theorem package proved here.

For the purposes of the present note, the important point
is only historical: the finite table came first. The
Fourier and $L$-value interpretation was discovered later.

% ============================================================
\section{Remarks}

\subsection*{Reading the table}

The collision periodic table says: every prime carries a
collision fingerprint determined by its residue class
modulo~$b^2$. In base~$10$, the prime $109$ and the
prime $50009$ have the same fingerprint ($S = +8$) because
both end in $09$. The prime $191$ and the prime $99991$
share $S = -9$ because both end in $91$. The structure is
periodic, exact, and computable.

\subsection*{The $-1$}

The antisymmetry constant is $-1$, not $0$. The offset
arises from the bilateral parity of the collision count
($C$ is even) combined with the relationship between
complementary bin structures. The complement map
$r \mapsto p - r$ sends digit~$d$ to digit~$b{-}1{-}d$,
exchanging large and small bins. The $-1$ accounts for the
parity mismatch between $Q$ and $Q'$ in paired classes.

\subsection*{The negative bias}

Exactly half the classes are negative. This is a structural
property of the antisymmetry: since $\mathcal{T}(a) +
\mathcal{T}(-a) = -1 < 0$, every pair contains at least one
negative member. The negative bias in the raw collision
fluctuation is explained by this: the collision invariant at
the specific multiplier $b$ is systematically below the mean, and the
periodic table shows this is not a statistical fluctuation
but a structural certainty.

% ============================================================
\begin{thebibliography}{9}

\bibitem{nfield}
A.~S.~Petty,
\emph{nfield: A structural analysis engine for fractional
field invariants}, software.
\url{https://github.com/alexspetty/nfield}

\end{thebibliography}

\end{document}
